% Generated by paperkit. Do not edit.
\newcommand{\PaperTitle}{When Should a Shopping Agent Stop Searching?}
\newcommand{\PaperSubtitle}{A measured decision rule for adaptive stopping under hard budgets}
\newcommand{\PaperAuthors}{Ahnaf Prio\\{\small Independent Researcher}}
\newcommand{\PaperAbstract}{An agent that shops on a user's behalf must decide, after every merchant it queries, whether to buy what it found or keep looking. Engineers usually pick a fixed rule, and nobody knows when that is good enough. We measure the question on live merchant data rather than argue it from theory. Two dated snapshots of a large set of protocol-enabled merchants show that offers really are ephemeral: listings are delisted and prices move within a month, while repeated scans minutes apart are byte-identical, so the measurement itself adds no noise. Replaying eight policies against frozen held-out panels built from these snapshots, an adaptive reservation-price rule attains lower failure-penalized purchase loss than the strongest fixed rule tuned on calibration data, with the whole paired bootstrap interval below zero and no budget violations. The advantage is statistically real but economically small, because cross-merchant price dispersion in this corpus is usually negligible and almost every product is carried by only two merchants. A sweep calibrated to the measured churn and dispersion maps where the rule pays for itself and where it does not: adaptive stopping helps once prices differ across merchants and the budget allows acting on what a query reveals, is indistinguishable from a tuned fixed rule when prices are nearly identical, and can be worse when the budget is too tight to use the information. We release the optimizer, the replay harness, and the collected panels.}
\newcommand{\PaperKeywords}{autonomous shopping, online search, optimal stopping, resource constraints, ephemeral offers}
\newcommand{\PaperArxivCategories}{cs.AI}
